\begin{lemma}
\label{lemma-artinian-finite-length}
A ring $R$ is Artinian if and only if it has finite length
as a module over itself. Any such ring $R$ is both Artinian and
Noetherian, any prime ideal of $R$ is a maximal ideal, and $R$ is equal
to the (finite) product of its localizations at its maximal ideals.
\end{lemma}

\begin{proof}
If $R$ has finite length over itself then it satisfies both
the ascending chain condition and the descending chain
condition for ideals. Hence it is both Noetherian and Artinian.
Any Artinian ring is equal to product of its localizations
at maximal ideals by Lemmas \ref{lemma-artinian-finite-nr-max},
\ref{lemma-artinian-radical-nilpotent}, and \ref{lemma-product-local}.

\medskip\noindent
Suppose that $R$ is Artinian. We will show $R$ has finite
length over itself. It suffices to exhibit a chain of
submodules whose successive quotients have finite length.
By what we said above
we may assume that $R$ is local, with maximal ideal $\mathfrak m$.
By Lemma \ref{lemma-artinian-radical-nilpotent} we have
$\mathfrak m^n =0$ for some $n$.
Consider the sequence
$0 = \mathfrak m^n \subset \mathfrak m^{n-1} \subset
\ldots \subset \mathfrak m \subset R$. By Lemma
\ref{lemma-dimension-is-length} the length of each subquotient
$\mathfrak m^j/\mathfrak m^{j + 1}$ is the dimension of this
as a vector space over $\kappa(\mathfrak m)$. This has to be
finite since otherwise we would have an infinite descending
chain of sub vector spaces which would correspond to an
infinite descending chain of ideals in $R$.
\end{proof}